\subsection{Transposition}

Let E be a locally convex space and \(E^{\prime}\) its dual. Recall that we denote by \(E_{b}^{\prime}\) and \(E_{c}^{\prime}\) the locally convex spaces obtained by endowing the vector space \(E^{\prime}\) with the topology of bounded convergence and that of compact convergence respectively (EVT, III, p. 14). The space \(E_{b}^{\prime}\) is also called the strong dual of E (loc. cit., example 4).

PROPOSITION 9. --- Let E be a locally convex space, F a separated locally convex space, and u a continuous linear map from E into F.

\begin{enumerate}
\def\labelenumi{\alph{enumi})}
\item
  If the mapping \(u\) is compact, then its transpose \(^{t}u\) is a compact linear map from \(F_c'\) into \(E_c'\) and a continuous linear map from \(F_c'\) into \(E_b'\) ;
\item
  Suppose that the space E is a semi-normed space and that the space F is quasi-complete. If the transpose \({}^{t}u\) is a continuous map from \(F_{c}^{\prime}\) into \(E_{b}^{\prime}\) then the linear map u is compact.
\end{enumerate}

Suppose first that the linear map \(u\) is compact. There then exists a neighborhood V of 0 in E whose image under \(u\) is contained in a compact subset C of F. By its definition (EVT, II, p. 47, def. 2 and p. 68, def. 1), the polar \(\mathrm{C}^{\circ}\) of C is a neighborhood of 0 in \(\mathrm{F}_c^\prime\) and we have \({}^t u(\mathrm{C}^\circ) \subset \mathrm{V}^\circ\) . Now \(\mathrm{V}^\circ\) is an equicontinuous subset of \(\mathrm{E}'\) closed for the weak topology, hence compact in \(\mathrm{E}_c^\prime\) (EVT, III, p. 17, cor. 2 and prop. 5). This proves that \({}^t u\) is a compact linear map from \(\mathrm{F}_c^\prime\) into \(\mathrm{E}_c^\prime\) .

Let us keep the preceding hypotheses and let U be a neighborhood of 0 in \(E_{b}^{\prime}\) . It contains the polar \(A^{\circ}\) of a bounded subset A of E. Since the linear map u is compact, the set \(u(\mathrm{A})\) is relatively compact in F (remark 1 of III, p. 2), and \((^{t}u)^{-1}(\mathrm{A}^{\circ}) = u(\mathrm{A})^{\circ}\) is a neighborhood of 0 in \(F_{c}^{\prime}\) . This proves that \(^{t}u\) is a continuous linear map from \(F_{c}^{\prime}\) into \(E_{b}^{\prime}\) .

Let us now place ourselves under the hypotheses of \(b\) and denote by B the unit ball of E. The polar \(\mathrm{B}^{\circ}\) of B is the unit ball of \(\mathrm{E}_b'\) and we have \((^{t}u)^{-1}(\mathrm{B}^{\circ}) = u(\mathrm{B})^{\circ}\) . If \(^{t}u\) is a continuous map from \(\mathrm{F}_c'\) into \(\mathrm{E}_b'\) , the set \(u(\mathrm{B})^{\circ}\) is therefore a neighborhood of 0 in \(\mathrm{F}_c'\) ; it then contains the polar \(\mathrm{C}^{\circ}\) of a compact subset C of F. By the bipolar theorem (EVT, II, p. 49, cor. 3), the set \(u(\mathrm{B})\) is contained in the closed convex hull of \(\mathrm{C} \cup \{0\}\) ; the latter is compact because the space F is quasi-complete (EVT, III, p. 8). This proves that the linear map \(u\) is compact.

COROLLARY 1 (Schauder). --- Let E be a seminormed space, F a Banach space, E' and F' their respective strong duals, and u a continuous linear map from E into F. The following properties are equivalent:

\begin{enumerate}
\def\labelenumi{(\roman{enumi})}
\item
  The linear map u from E into F is compact;
\item
  The linear map \({}^{t}u\) from \(\mathrm{F}'\) into \(\mathrm{E}'\) is compact;
\item
  The linear map \({}^{t}u\) from \(\mathrm{F}_{c}^{\prime}\) into \(\mathrm{E}'\) is continuous.
\end{enumerate}

The equivalence of (i) and (iii) follows from Prop. 9, a) and b); the implication (iii)⇒(ii) comes from the fact that the identity map from F' into F'c is compact (cor. of Prop. 7 of III, p. 7).

Let us conclude by proving that condition (ii) implies (i). Suppose that the linear map \({}^{t}u\) : \(F^{\prime} \to E^{\prime}\) is compact. It follows from the implication (i)⇒(ii), applied to \({}^{t}u\) , that \({}^{tt}u\) is a compact linear map from \(E''\) into \(F''\) . Denote by \(v\) the canonical linear map from E into \(E''\) ; it is continuous. Since F is identified with a closed vector subspace of \(F''\) and u coincides with \({}^{t}({}^{t}u) \circ v\) on E, it follows from Remark 5 of III, p. 3 that the linear map \(u\) is compact.

COROLLARY 2. --- Let E be a semi-normed space and F a Banach space of countable type. Let u be a continuous linear map from E into F. Suppose that for every sequence \((y_{n})\) of elements of F' weakly tending to 0, the sequence \((^{t}u(y_{n}))\) tends strongly to 0 in E'. Then the linear map u is compact.

Let B′ denote the unit ball of F′, equipped with the topology induced by the topology \(\sigma(\mathrm{F}',\mathrm{F})\) . Since F is a Banach space of countable type, B' is a compact metrizable space (EVT, III, p. 19, cor. 2). Under the hypothesis made, the restriction of \({}^{t}u\) to B' is a continuous map from B' into E' and \({}^{t}u(\mathrm{B}')\) is a compact subset of E'. By Cor. 1, the linear map u is compact.

Remark. --- Let E and F be separated locally convex spaces. The transpose of a compact linear map from E into F is not always a compact linear map from \(F_{b}^{\prime}\) into \(E_{b}^{\prime}\) (cf. III, p. 108, exercise 15).

